% ===========================================================================
\section{Doğrulama}\label{sec:dogrulama-aile}
% ===========================================================================

Donanımla ilgili her iddia, tamsayı referansla bit düzeyinde eşitliğe dayanır. Tablo~\ref{tab:dogrulama-aile} bu
eşitliğin hangi yollarla gösterildiğini özetler. Bütün Verilator simülasyonları doğrulamalar (assertion) açık olarak
çalışır. Bu doğrulamalar şunları kontrol eder:
\begin{itemize}
  \item bellek adreslerinin aralık içinde kalması,
  \item INT32 akümülatörünün hiçbir şeritte taşmaması,
  \item bağlam konumunun $T$'den küçük olması,
  \item sınırlı bölümlü bölücünün ön koşulu ($\mathrm{num}\gg Q_B < \mathrm{den}$).
\end{itemize}

\begin{table}[h]
  \centering\small
  \caption[Doğrulama kanıtları]{Doğrulama kanıtları. Bütün satırlarda hata sayısı sıfırdır.}
  \label{tab:dogrulama-aile}
  \begin{tabularx}{\textwidth}{@{}lL@{}}
    \toprule
    \thead{Kontrol} & \thead{Kapsam} \\ \midrule
    Eğitilmiş her model (pipeline) & \zooVerifRuns{} model; test kümesinin tamamı, \zooVerifiedLogits{} logit \\
    Özelleşmiş motor, yükleme portu & aynı test vektörleri \\
    Tek bitstream motoru, AXI üzerinden & 10 model; 200 tokenlık kapalı döngü hikâye, token token aynı \\
    Rastgele farklılık testi (\code{fuzz.py}) & \zooFuzzCases{} rastgele mimari/boyut/ağırlık, \zooFuzzLogits{} logit \\
    4 durumlu simülasyon (Icarus) & çıktılarda belirsiz (X) değer yok \\
    Aritmetik birimler (\code{tb\_zarith.sv}) & 12 bitlik bölücüde bütün geçerli girişler; toplam \zooArithCases{} vaka \\
    Bağımsız C ISA simülatörü & 10 modelin altın vektörleri \\
    Sürekli entegrasyon (GitHub Actions) & her commit'te 9 model + C simülatörü + AXI sistemi \\
    Kart üzerinde (PYNQ-Z2) & \zooBoardVerif{} \\ \bottomrule
  \end{tabularx}
\end{table}

\paragraph{Rastgele farklılık testi.} Eğitilmiş modeller, ağırlık uzayının yalnızca küçük bir bölgesini dener. Fuzzing
her vakada rastgele seçimler yapar: bir mimari, sözlük boyutu (27--512), bağlam (8--64), derinlik (1--3) ve genişlik.
Ağırlıklar rastgele ölçekli değerlerle doldurulur; böylece doyma ve kırpma yolları da sınanır. Model rastgele tokenlarla
kalibre edilir ve rastgele uzunluktaki rastgele dizilerin her logiti karşılaştırılır. Derleyicinin donanım biçimlerine
sığmadığı için reddettiği vakalar ayrıca sayılır. İlk sürümde bir vaka böyle reddedildi: artık akış çok küçük
değerlerdeyken derleyici, normalize vektörün INT16'ya sığmayacağı bir ölçek seçiyordu. Derleyici bu durum için
düzeltildi.

\paragraph{Aritmetik birimler.} Sınırlı bölümlü bölücünün doğruluğu iki yoldan gösterilir. (i)~$W=12$, $Q_B=6$ için ön
koşulu sağlayan \emph{bütün} giriş çiftleri simüle edilir; algoritma genişlikten bağımsız olduğu için bu bir kanıt
yerine geçer. (ii)~Gerçek genişliklerde rastgele ve sınır vakalar denenir. Softmax'ın karşılığı ise $[2^{15}, 2^{24})$
aralığındaki her payda için kontrol edilir. Sonuçlar ikinci bir uygulamayla değil, tanımlayıcı özelliklerle denetlenir:
$q\,b\le a<(q+1)\,b$ ve $r^2\le n<(r+1)^2$.

\paragraph{Üç bağımsız uygulama.} Tamsayı semantiğinin üç ayrı uygulaması vardır:
\begin{itemize}
  \item Python referansı (\code{IntRunner}),
  \item RTL,
  \item mikrokodu yorumlayan bir C programı (\code{sw/c/zoo\_run.c}).
\end{itemize}
C programı aynı bellek görüntülerini okur ve on modelin hepsinde altın vektörlerle bit düzeyinde aynıdır. Kartta ARM
çekirdeğinde de aynı program çalışır; bu yüzden CPU karşılaştırması, birebir aynı tamsayı işinin karşılaştırmasıdır.

\paragraph{Sentez sonrası.} Yosys'in Xilinx hücre modelleriyle yapılan netlist simülasyonu, bu modellerin ilk değeri
atanmamış iç yazmaçlarından kaynaklanan belirsiz (X) değerler nedeniyle sonuç vermedi. Gerçek FPGA'de yazmaçlar sıfırla
başlar ve RTL düzeyindeki 4 durumlu simülasyon temizdir. Sentez sonrası kontrol, Vivado'nun hücre modelleriyle yapılacak
şekilde betik olarak hazırlanmıştır (\code{scripts/vivado\_zoo\_postsynth\_sim.tcl}).
